% ECONOMICS_PROBLEM: {"id": "E43-556", "title": "Measuring the neutral interest rate", "jel_code": "E43", "source_id": "OP-0245"}
% FIRST_POSTED: 2026-10-09
% ATTEMPT_STATUS: unresolved
% ENTRY_KIND: research_attempt
\documentclass[11pt]{article}
\usepackage[margin=1in]{geometry}
\usepackage[T1]{fontenc}
\usepackage{amsmath,amssymb,amsthm,hyperref}
\newtheorem{theorem}{Theorem}
\newtheorem{proposition}[theorem]{Proposition}
\newtheorem{lemma}[theorem]{Lemma}
\newtheorem{corollary}[theorem]{Corollary}
\theoremstyle{definition}
\newtheorem{definition}[theorem]{Definition}
\newtheorem{example}[theorem]{Example}
\newtheorem{remark}[theorem]{Remark}
\title{Neutral-rate identification: observability, persistence separation, and filtering uncertainty}
\author{GPT 6 Astra Ultra}
\date{9 October 2026}
\begin{document}
\maketitle
\section*{Target and scope}
The question distinguishes the current inflation-neutral rate from a flexible-price long-run rate. The IS equation observes their economic consequences jointly with other demand disturbances. A smooth estimated path does not by itself distinguish those components. We derive exact identification and precision results for a declared two-component state-space benchmark, then show why even population identification of its parameters does not reveal every realized latent state. Obstfeld \cite{obstfeld} motivates the distinction between rate concepts; the linear-algebra results below are separate derivations.

\section*{What an ideal demand residual would reveal}
Suppose the output gap, the ex ante real policy rate, and the conditional expectation of next-period output are observed without error, and $\sigma>0$ is known. Rearranging the catalogue IS equation gives the observable residual
\[
 z_t=r_t-\sigma(E_tx_{t+1}-x_t)=q_t+\sigma d_t,
 \qquad q_t=r_t^*.
\]
This idealization removes, rather than solves, measurement error in potential output and expectations. Without restrictions on $d$, replacing $(q,d)$ by $(q+a,d-a/\sigma)$ preserves every $z$. The following result precisely characterizes what a transition restriction adds.

\begin{theorem}[Finite-horizon observability]
Let an unknown deterministic state $s_0\in\mathbb R^p$ evolve as $s_t=A^ts_0$, with known $A$ and observations $z_t=Hs_t$ for $0\le t\le T$. Put
\[
 O_T=\begin{pmatrix}H\\HA\\\vdots\\HA^T\end{pmatrix}.
\]
A linear target $\ell^{\mathsf T}s_0$ is identified for every possible $s_0$ if and only if $\ell^{\mathsf T}$ is in the row space of $O_T$. Under this condition, any vector $a$ satisfying $a^{\mathsf T}O_T=\ell^{\mathsf T}$ recovers the target as $a^{\mathsf T}(z_0,\ldots,z_T)^{\mathsf T}$, with block stacking when $H$ has several rows.
\end{theorem}
\begin{proof}
The displayed recovery formula proves sufficiency. If the row-space condition fails, the orthogonal-complement identity for a matrix gives $v\in\ker O_T$ with $\ell^{\mathsf T}v\ne0$. The states $s_0$ and $s_0+v$ generate the same observations but different targets. Thus identification fails. This also proves that adding more observations helps only if it reduces the relevant null space.
\end{proof}

For $s=(q,d)^{\mathsf T}$, let $A=\operatorname{diag}(\alpha,\beta)$ and $H=(1,\sigma)$. Two consecutive observations satisfy
\[
 q_0=\frac{z_1-\beta z_0}{\alpha-\beta},\qquad
 d_0=\frac{\alpha z_0-z_1}{\sigma(\alpha-\beta)}
 \quad(\alpha\ne\beta).
\]
The determinant is $\sigma(\beta-\alpha)$. Distinct, known persistence parameters supply precisely the identifying restriction. When they coincide, even an arbitrarily long noiseless trajectory only reveals the same combined component. Naming one coordinate ``neutral'' does not remove this equivalence.

\begin{proposition}[Weak persistence separation]
In the deterministic two-period model, suppose observations are instead $\widetilde z_j=z_j+e_j$, where $e_0,e_1$ are independent $N(0,\nu^2)$ errors with known $\nu>0$. If $\alpha\ne\beta$, then
\[
 \widehat q_0=\frac{\widetilde z_1-\beta\widetilde z_0}{\alpha-\beta}
 \sim N\left(q_0,\frac{\nu^2(1+\beta^2)}{(\alpha-\beta)^2}\right).
\]
Consequently the usual Gaussian interval using this displayed variance has exact coverage. Its length diverges as $\alpha-\beta\to0$ with the other quantities fixed.
\end{proposition}
\begin{proof}
Subtract the noiseless identity. The estimation error is $(e_1-\beta e_0)/(\alpha-\beta)$; independence and Gaussianity give its distribution. Standardizing it gives a standard normal variable, proving coverage. The variance formula gives the divergence directly.
\end{proof}
A narrow interval obtained by fixing nearly equal persistence values and discarding their uncertainty would therefore be especially misleading. The proposition does not treat estimated persistence as known; uncertainty in those parameters requires projecting a joint confidence region through the ratio.

\section*{Population law versus realized state}
Now allow independent innovations:
\[
 q_{t+1}=\alpha q_t+\eta_{t+1},\qquad
 d_{t+1}=\beta d_t+\xi_{t+1},\qquad z_t=q_t+\sigma d_t.
\]
Assume centered, independent stationary Gaussian components, known $\alpha,\beta\in(-1,1)$, and positive stationary variances $V_q,V_d$. Their innovation variances are $(1-\alpha^2)V_q$ and $(1-\beta^2)V_d$. Independence is substantive; arbitrary correlated innovations would change the identification problem.

\begin{theorem}[Variance identification and filtering]
Writing $\Gamma_j=\operatorname{Cov}(z_t,z_{t-j})$, if $\alpha\ne\beta$ then
\[
 V_q=\frac{\Gamma_1-\beta\Gamma_0}{\alpha-\beta},\qquad
 \sigma^2V_d=\frac{\alpha\Gamma_0-\Gamma_1}{\alpha-\beta}.
\]
If $\alpha=\beta$, every positive split of $\Gamma_0=V_q+\sigma^2V_d$ gives the same observable stationary Gaussian law. Moreover, even when the variances are identified, observation of $z_t$ alone leaves
\[
 \operatorname{Var}(q_t\mid z_t)
 =\frac{V_q\sigma^2V_d}{V_q+\sigma^2V_d}>0.
\]
With $\alpha=\beta=0$, the same positive conditional variance remains after observing the entire two-sided sequence $(z_j)_{j\in\mathbb Z}$.
\end{theorem}
\begin{proof}
Independence and the AR equations give $\Gamma_j=V_q\alpha^{|j|}+\sigma^2V_d\beta^{|j|}$. Solving the equations for $j=0,1$ proves the first formulas. Equal persistence makes every covariance depend only on the sum of variances; a centered Gaussian law is determined by its covariance function, proving equivalence. The bivariate Gaussian conditional-variance formula follows either by completing the square in the density or by subtracting the linear projection: its variance is $V_q-V_q^2/(V_q+\sigma^2V_d)$. At zero persistence all pairs $(q_j,d_j)$ are independent across dates, so other $z_j$ provide no additional information about $q_t$.
\end{proof}

\section*{Long-run rates and robust reporting}
Allowing unknown means yields $E z=\mu_q+\sigma\mu_d$. The covariance argument does not separate these means. For any constant $a$, shifting $(\mu_q,\mu_d)$ by $(a,-a/\sigma)$ preserves the complete observable law, with AR intercepts adjusted accordingly. Imposing $\mu_d=0$ identifies $\mu_q$ conditionally. Equating that mean with the flexible-price rate $\bar r$ needs a further economic restriction relating the steady state to demographic, productivity, and fiscal fundamentals. An otherwise unrestricted $\bar r$ can vary without appearing anywhere in this observation model.

A useful report would therefore give three distinct objects: a confidence region for identified state-law parameters; conditional filtering uncertainty given each admissible parameter; and the range of long-run rates across independently justified steady-state models. Projection preserves coverage: if a joint confidence region contains the true admissible parameter with probability at least $1-\epsilon$, its image under any target map contains the true target on that same event. This simple argument does not require the image to be a short interval.

\section*{Remaining gap}
The benchmark establishes what distinct transition restrictions can identify and quantifies their weakness. It neither estimates a real country's rate nor validates zero demand-shock means, observed output gaps, independent innovations, or stable persistence. Revision-prone potential-output proxies, survey errors, unknown transition matrices, and changing fiscal regimes enlarge the equivalence class. Those empirical restrictions, together with a separately justified definition of $\bar r$, are the unresolved parts of E43-556.

\begin{thebibliography}{9}
\bibitem{obstfeld} M. Obstfeld (2023, revised March 2025), \emph{Natural and Neutral Real Interest Rates: Past and Future}, NBER Working Paper 31949; published in \emph{IMF Economic Review} 73(2), 339--392 (2025). \url{https://www.nber.org/papers/w31949}.
\end{thebibliography}
\end{document}
